\documentclass[10pt,twocolumn,letterpaper]{article}

\usepackage[top=0.75in,bottom=0.75in,left=0.75in,right=0.75in]{geometry}
\usepackage{times}
\usepackage{amsmath}
\usepackage{amssymb}
\usepackage{booktabs}
\usepackage{graphicx}
\usepackage{hyperref}
\usepackage{microtype}
\usepackage{caption}
\usepackage{subcaption}

\hypersetup{
    colorlinks=true,
    linkcolor=blue,
    filecolor=magenta,      
    urlcolor=cyan,
    pdftitle={COS30082 Assignment 1: Bird Species Classification},
}

\title{\textbf{Fine-Grained Bird Species Classification on CUB-200 via Modernized Pure Convolutions and Regularized Transfer Learning}}
\author{
    \textbf{Andy Wei Jian CHU} \\
    Student ID: \texttt{104401283} \\
    Department of Computing Technologies, Swinburne University of Technology \\
    COS30082 Applied Machine Learning --- Assignment 1
}
\date{October 2026}

\begin{document}

\maketitle

\begin{abstract}
Fine-Grained Visual Categorization (FGVC) presents a formidable challenge in computer vision due to extremely low inter-class morphological variance across visually proximate species, compounded by large intra-class pose, lighting, and background variability. This paper investigates bird species identification on the Caltech-UCSD Birds 200 (CUB-200) benchmark (200 bird categories, 4,829 training and 1,204 testing images). To prevent catastrophic overfitting on this low-sample dataset ($\approx 24$ images/class), we establish a stratified 85/15 validation regime, domain-specific data augmentations (Random Resized Crop, Color Jitter, Random Erasing), Label Smoothing Cross-Entropy, and AdamW optimization with Cosine Annealing. We contrast a canonical baseline (\textbf{ResNet-50}) against a modernized pure convolutional architecture (\textbf{ConvNeXt-Tiny}). Experimental evaluation reveals that ConvNeXt-Tiny decisively outperforms ResNet-50, achieving an overall \textbf{Top-1 Accuracy of 89.53\%} and an \textbf{Average Accuracy per Class of 89.36\%} (an absolute gain of $+8.88\%$ and $+9.25\%$ respectively). Finally, we conduct fine-grained confusion pair analyses and deploy the optimized model as an interactive web service via Gradio on Hugging Face Spaces.
\end{abstract}

\vspace{0.5em}
\noindent\textbf{Hugging Face Web Application:} \url{https://huggingface.co/spaces/AA-49/Bird-Species-Classification} \\

\section{Methodology}

\subsection{Problem Formulation and Dataset Complexity}
The Caltech-UCSD Birds 200 (CUB-200) dataset consists of 200 distinct bird species. We are provided with 4,829 training samples and 1,204 testing samples annotated with species labels $y \in \{0, 1, \dots, 199\}$. Fine-grained bird classification exhibits two distinct properties:
\begin{enumerate}
    \item \textbf{Subtle Inter-Class Discriminators:} Sibling species (e.g., \textit{Laysan Albatross} vs. \textit{Black-footed Albatross}, or sympatric sparrows and warblers) share virtually indistinguishable silhouettes, relying on regional attributes such as bill morphology, eye rings, or plumage barring.
    \item \textbf{Extreme Overfitting Risk:} Standard convolutional networks trained on small image counts readily memorize invariant background textures (e.g., foliage, open sky, branches) instead of bird morphology.
\end{enumerate}

\subsection{Architectural Paradigms}
To evaluate representation capacity, we compare ResNet-50 and ConvNeXt-Tiny under two regimes: a frozen feature extractor with only the new classifier head trained, and end-to-end fine-tuning with differential learning rates. This produces four configurations and enables frozen-vs-fine-tuned comparisons within each architecture and architecture comparisons at each regime. Existing production results below describe fine-tuned checkpoints; frozen-regime results must be added after the new experiment completes.

\subsubsection{ResNet-50}
ResNet-50 addresses the degradation problem through residual skip connections:
\begin{equation}
    \mathbf{y} = \mathcal{F}(\mathbf{x}, \{W_i\}) + \mathbf{x}
\end{equation}
The architecture organizes bottleneck blocks ($1\times1$, $3\times3$, $1\times1$ convolutions) across four stages, producing a 2048-dimensional feature vector via Global Average Pooling (GAP). We detach the ImageNet-1K classifier and append a custom classification head:
\begin{equation}
    \mathbf{z} = \mathbf{W}_{\text{head}} \cdot \text{Dropout}_{p=0.3}(\text{GAP}(\mathbf{F})) + \mathbf{b}_{\text{head}}
\end{equation}
where $\mathbf{W}_{\text{head}} \in \mathbb{R}^{200 \times 2048}$. While computationally reliable, ResNet's localized $3\times3$ spatial kernels limit receptive field expansion until the deepest layers, constraining its ability to jointly model distant diagnostic regions (e.g., associating bill color with tail plumage).

\subsubsection{ConvNeXt-Tiny}
ConvNeXt-Tiny modernizes convolutional neural networks by incorporating design principles established by the Swin Transformer:
\begin{itemize}
    \item \textbf{Inverted Bottleneck Design:} ConvNeXt inverts channel expansion ratios, adopting hidden dimensions 4 times wider than input dimensions. This preserves high-frequency feather texture signals.
    \item \textbf{Large $7\times7$ Depthwise Convolutions:} Replacing $3\times3$ kernels with $7\times7$ depthwise spatial convolutions expands the effective receptive field early in the network hierarchy, matching Vision Transformer attention spans while preserving translation equivariance.
    \item \textbf{Micro-Design Enhancements:} Replacing standard ReLU and BatchNorm layers with Gaussian Error Linear Units (GELU) and Layer Normalization (LN) substantially stabilizes training dynamics across fine-grained channel manifolds.
\end{itemize}

\subsection{Overfitting Mitigation and Regularization Suite}
To satisfy the rubric requirement of addressing overfitting, we implement a four-part regularization suite:

\subsubsection{Stratified 85/15 Data Partitioning}
To enable reliable early stopping and hyperparameter tuning without contaminating test results, the 4,829 training images were partitioned into 4,104 training and 725 validation images using stratified random sampling ($\text{seed}=42$), ensuring uniform representation across all 200 species.

\subsubsection{Domain-Aware Data Augmentation}
Input images are dynamically augmented during training:
\begin{itemize}
    \item \textbf{Random Resized Crop ($224\times224$, scale $0.5\text{--}1.0$):} Promotes scale invariance and forces the network to discover localized body parts.
    \item \textbf{Random Horizontal Flip ($p=0.5$):} Enforces bilateral pose symmetry.
    \item \textbf{Color Jitter:} Stochastically perturbs brightness ($\pm 20\%$), contrast ($\pm 20\%$), saturation ($\pm 20\%$), and hue ($\pm 5\%$) to neutralize outdoor daylight variations.
    \item \textbf{Random Erasing ($p=0.2$, scale $0.02\text{--}0.2$):} Obscures rectangular patches to simulate occluding foliage and tree branches.
\end{itemize}

\subsubsection{Label Smoothing Loss Function}
To mitigate overconfidence on visually ambiguous species, we replace hard categorical targets with \textbf{Label Smoothing Cross-Entropy} ($\epsilon = 0.1$):
\begin{equation}
    y_k^{\text{smooth}} = (1 - \epsilon)\delta_{k, y} + \frac{\epsilon}{C}
\end{equation}
\begin{equation}
    \mathcal{L}_{\text{LS}}(\mathbf{p}, \mathbf{y}) = - \sum_{k=1}^C y_k^{\text{smooth}} \log(p_k)
\end{equation}
where $C=200$. This prevents output logits from diverging to infinity, maintaining margin separability between visually adjacent bird species.

\subsubsection{Optimization and Mixed Precision Scheme}
Training is orchestrated using \textbf{AdamW} with decoupled weight decay ($\lambda = 0.01$). We employ differential learning rates: $\eta_{\text{backbone}} = 1\times 10^{-4}$ and $\eta_{\text{head}} = 5\times 10^{-4}$, scheduled via Cosine Annealing down to $\eta_{\text{min}} = 1\times 10^{-6}$. To fit within the 4GB VRAM constraint of the local NVIDIA GeForce RTX 2050 GPU, Automatic Mixed Precision (AMP FP16) was activated, cutting memory consumption to $\approx 1.8$\,GB.

\section{Result and Discussion}

\subsection{Evaluation Metrics and Test Results}
In strict adherence to the project brief, performance is reported using two required evaluation metrics on the 1,204 official test images:

\begin{equation}
    \text{Top-1 Accuracy} = \frac{1}{N} \sum_{k=1}^N \mathbf{1}\{\text{argmax}(y) == \text{ground\_truth}\}
\end{equation}

\begin{equation}
    \text{Average Accuracy per Class} = \frac{1}{C} \sum_{i=1}^C T_i
\end{equation}
where $T_i$ is the accuracy computed across all test samples belonging to class $C_i$, and $C = 200$.

\begin{table}[h]
\centering
\small
\caption{Quantitative Comparison on Official CUB-200 Test Set.}
\label{tab:results}
\begin{tabular}{lccc}
\toprule
\textbf{Metric} & \textbf{ResNet-50} & \textbf{ConvNeXt-Tiny} & \textbf{$\Delta$ Gain} \\
\midrule
\textbf{Top-1 Accuracy} & 80.65\% & \textbf{89.53\%} & \textbf{+8.88\%} \\
\textbf{Average Accuracy per Class} & 80.11\% & \textbf{89.36\%} & \textbf{+9.25\%} \\
Parameters & 24.5M & 28.0M & +14.2\% \\
Peak VRAM (FP16) & 1.12 GB & 1.18 GB & +5.3\% \\
Final Val Loss & 1.542 & \textbf{1.318} & -14.5\% \\
\bottomrule
\end{tabular}
\end{table}

\subsection{Architectural Comparative Analysis}
As delineated in Table~\ref{tab:results}, \textbf{ConvNeXt-Tiny decisively outperforms ResNet-50 across both official metrics}, securing $+8.88\%$ in Top-1 Accuracy ($89.53\%$ vs. $80.65\%$) and $+9.25\%$ in Average Class Accuracy ($89.36\%$ vs. $80.11\%$).

The current baseline's ConvNeXt-Tiny result is consistent with three architectural hypotheses:
\begin{enumerate}
    \item \textbf{Receptive Field Dynamics:} ConvNeXt's $7\times7$ depthwise kernels correlate spatial landmarks (such as head markings with wing bars) at earlier layers. ResNet-50 relies on stacked $3\times3$ convolutions which diffuse local details before global contextual integration occurs.
    \item \textbf{Inverted Bottleneck Spatial Preservation:} By expanding channel width in intermediate block stages, ConvNeXt avoids early dimension bottlenecks, preserving high-frequency feather texture representations.
    \item \textbf{Optimization Landscape Stability:} Layer Normalization provides consistent normalization across channel axes, suppressing internal covariate shift without dependence on mini-batch statistics.
\end{enumerate}

\subsection{Fine-Grained Error Analysis and Confusion Insights}
The saved baseline evaluation reports the following directed confusions. They are qualitative review candidates, not evidence that confusion is strictly taxonomic:
\begin{itemize}
    \item \textbf{Common Tern (class 143) predicted as Arctic Tern (class 140):} 3 test errors in this direction.
    \item \textbf{Chuck-will's-widow (class 21) predicted as Whip-poor-will (class 104):} 2 errors in this direction; the reverse direction also has 2.
    \item \textbf{Acadian Flycatcher (class 36) predicted as Least Flycatcher (class 38):} 2 errors in this direction; the reverse direction also has 2.
\end{itemize}

\subsection{Overfitting Diagnosis}
The saved fine-tuned histories show a final training-to-validation Top-1 gap of 15.83 percentage points for ResNet-50 (99.42% vs. 83.59%) and 11.20 points for ConvNeXt-Tiny (99.61% vs. 88.41%). These histories document runs with the stated regularization configuration; no no-augmentation/no-smoothing ablation artifacts are present to establish causal effects or support a below-7% gap claim. Frozen-versus-fine-tuned runs compare training regimes, not individual regularization components.

\section{Model Deployment and Web Application}
The current baseline checkpoints are encapsulated into an interactive web application utilizing \textbf{Gradio} and prepared for \textbf{Hugging Face Spaces}. A new checkpoint will only be promoted after the 12-run comparison selects it by validation average class accuracy and its official test evaluation is saved. The user interface provides:
\begin{itemize}
    \item Multi-source image input (upload, clipboard, webcam).
    \item Top-5 predicted species probability distribution with confidence metrics and clean taxonomical common names.
    \item Bundled benchmark examples from the official test dataset for instant one-click demonstration.
\end{itemize}

\section{Conclusion}
The current baseline achieves \textbf{89.53\% Top-1 Accuracy} and \textbf{89.36\% Average Class Accuracy} on CUB-200. The frozen-versus-fine-tuned comparison remains pending; production promotion requires validation-based selection followed by official test evaluation.

\end{document}
